% %%%%%%%%%%%%%%%%%%%%%%%%%%%%%%%%%%%%%%%%%%%%%%%%%%%%%%%%%%%%%%%%%%%%%%%%%%%%
% Thesis Defence
% Section: State-of-the-Art
% %%%%%%%%%%%%%%%%%%%%%%%%%%%%%%%%%%%%%%%%%%%%%%%%%%%%%%%%%%%%%%%%%%%%%%%%%%%%

\section[\'Etat-de-l'art]{Accélérateurs dédiées et directions}

\subsection{Accélérateurs matériels dédiés}

\refStep{ref:GPU}

\begin{frame}
    \frametitle{Accélérateurs matériels dédiés}

    \pause[1]%
    {%
    \def\X{6cm}%
    \def\Y{2.35cm}%
    \begin{tikzpicture}[remember picture, overlay]
        \node [shift={(\X, \Y)}, single arrow, draw=Red!70!white, very thick, fill=Red!70!white, opacity=1., minimum width = 1.5cm, single arrow head extend=2pt, minimum height=18cm, rotate=0, blur shadow={shadow xshift=1pt,shadow yshift=-2pt}] at (current page.west)
            {}; % length of arrow
    \end{tikzpicture}%
    %
    \begin{tikzpicture}[remember picture, overlay]
        \node [shift={(\X - 5.1cm, \Y + 0.25cm)}, rotate=0] at (current page.west)
            {%
                \bfseries 1973%
            };
    \end{tikzpicture}%
    %
    \begin{tikzpicture}[remember picture, overlay]
        \node [shift={(\X - 0.95cm, \Y + 0.25cm)}, rotate=0] at (current page.west)
            {%
                \bfseries 2005%
            };
    \end{tikzpicture}%
    %
    \begin{tikzpicture}[remember picture, overlay]
        \node [shift={(\X + 3.75cm, \Y + 0.25cm)}, rotate=0] at (current page.west)
            {%
                \bfseries 2018%
            };
    \end{tikzpicture}%
    %
    \begin{tikzpicture}[remember picture, overlay]
        \node [shift={(\X + 7.4cm, \Y + 0.25cm)}, rotate=0] at (current page.west)
            {%
                \bfseries Aujourd'hui%
            };
    \end{tikzpicture}%
    %
    \begin{tikzpicture}[remember picture, overlay]
        \node [shift={(\X - 3.3cm, \Y - 0.5cm)}, rotate=0] at (current page.west)
            {%
                \lbox{black}{%
                    \parbox[t]{4cm}{\centering Approches logicielles sur CPUs \cref{ref:spFormats1}\vphantom{Ay}}
                }%
            };
    \end{tikzpicture}%
    %
    \begin{tikzpicture}[remember picture, overlay]
        \node [shift={(\X + 1.4cm, \Y - 0.5cm)}, rotate=0] at (current page.west)
            {%
                \lbox{black}{%
                    \parbox[t]{4cm}{\centering Approches logicielles pour GPUs \cref{ref:GPU}\vphantom{Ay}}
                }%
            };
    \end{tikzpicture}%
    %
    \begin{tikzpicture}[remember picture, overlay]
        \node [shift={(\X + 6.1cm, \Y - 0.5cm)}, rotate=0] at (current page.west)
            {%
                \lbox{black}{%
                    \parbox[t]{4cm}{\centering\bfseries Accélérateurs matériels dédiés au calcul creux\vphantom{Ay}}
                }%
            };
    \end{tikzpicture}%
    %
    % \begin{tikzpicture}[remember picture, overlay]
    %     \node [shift={(\X + 6.0cm, \Y - 1.2cm)}, rotate=0] at (current page.west) (A)
    %         {%
    %             \itshape\small Fin de la loi de Dennard
    %         };%
    %     \draw [Stealth-, draw=black, thick] (A.west) -- ++(-0.5,0) -- ++(0,1.2);
    % \end{tikzpicture}
    }

    \bigskip\bigskip\bigskip\smallskip

    \pause[2]%
    \textbf{\textcolor{Red}{Plusieurs caractéristiques}}

    \smallskip

    \point[\bfseries][2][Red]{\textit{Kernel} :} {\def\c{Orange} \lbox{\c}{Matrice-Matrice} \lbox{\c}{Matrice-Vecteur} \lbox{\c}{Générique}}

    \smallskip

    \point[\bfseries][2][Red]{Algorithme :} {\def\c{Green} \lbox{\c}{Inner-Product} \lbox{\c}{Outer-Product} \lbox{\c}{Gustavson} \lbox{\c}{Méthodes de \textit{tiling}}}

    \smallskip

    \point[\bfseries][2][Red]{Problèmes adressés :} {\def\c{Blue} \lbox{\c}{Gather (Intersections)} \lbox{\c}{Scatter (Réductions)}}

    \smallskip

    \point[\bfseries][2][Red]{Type d'accélérateur :} {\def\c{Cyan} \lbox{\c}{Autonome} \lbox{\c}{Assisté par processeur}}

    \smallskip

    \point[][2][Red]{Pour les accélérateurs assistés par processeur,} \textbf{position dans la hiérarchie mémoire :}

    {\def\c{Purple}\hspace*{0.6cm} \lbox{\c}{Proche processeur} \lbox{\c}{Proche caches} \lbox{\c}{\textit{Bypass} caches} \lbox{\c}{Proche mémoire principale}}

    \vspace{0.4cm}
    \drefX{ref:GPU}{Thaha Mohammed and Rashid Mehmood. 2022. Performance Enhancement Strategies for Sparse Matrix-Vector Multiplication (SpMV) and Iterative Linear Solvers. (2022), 41.}

\end{frame}

\begin{frame}
    \frametitle{Diversité des accélérateurs existants}
    
    \pause[1]%
    Sous-ensemble sur \textbf{\textcolor{Red}{43 accélérateurs}} étudiés entre 2015 et 2025.

    \vspace*{-5pt}

    \begin{minipage}[t]{0.68\linewidth}
        \centering

        \pause[1]%
        {
            \def\ck{\ding{51}}
            % \def\sa{\ding{86}\textsuperscript{1}}
            % \def\sb{\ding{86}\textsuperscript{2}}
            % \def\sc{\ding{86}\textsuperscript{3}}
            % \def\sd{\ding{86}\textsuperscript{4}}
            \def\sa{}
            \def\sb{$\mathbf{\sim}$}
            \def\sc{$\mathbf{\sim}$}
            \def\sd{$\mathbf{\sim}$}
            \begin{tblr}{
                    width = {\linewidth},
                    colspec = {X[2.5,c] X[1,c] *{12}{X[1,c]}}, % rowspec
                    rows = {m, font=\scriptsize},
                    colsep = {1pt},
                    row{1-Z} = {rowsep=0.5pt, black!2},
                    row{1-2} = {c, font=\bfseries\scriptsize, black!5},
                    cell{2}{3-14} = {font=\bfseries\itshape\scriptsize},
                    % row{W-Z} = {l, font=\itshape\scriptsize, rowsep=-2pt},
                    row{1,3} = {abovesep=1.5pt},
                    row{2,Z} = {belowsep=1.5pt},
                    cell{1}{1,2} = {r=2}{valign=f},
                    cell{1}{2} = {cmd=\rotatebox{90}},
                    cell{2}{3-14} = {cmd=\rotatebox{90}},
                    cell{1}{3,11,13} = {c=2}{},
                    cell{1}{5,8} = {c=3}{},
                    % cell{W-Z}{1} = {c=14}{},
                    hspan = {minimal},
                    stretch = 0,
                    % hlines, vlines,
                    % rulesep = {-1pt},
                    hline{1,Z} = {2pt}, % Z for last
                    hline{3} = {1pt},
                    % hline{Z} = {1pt, dashed},
                    vline{2,3,5,8,11,13} = {1-Z}{0.5pt},
                }
                %
                Nom & Année & Autonomie & & {Position dans la hiérarchie mémoire} & & & {Algorithme(s) le(s) plus adapté(s)} & & & {Support matériel pour} & & {\textit{Kernels} creux adressés} & \\
                & & Auton. & Assisté & Proc. & Caches & Mém. & IP & OP & Gust & Gather & Scatter & MV & MM \\
                %
                Coup                     & 2015 &     & \ck &     & \ck &     &     & \ck &     &     & \ck & \ck & \sc \\
                EIE                      & 2016 & \ck &     &     &     & \sa &     & \ck &     & \ck & \ck & \ck &     \\
                OuterSPACE               & 2018 & \ck &     &     &     & \sa &     & \ck &     & \ck & \ck & \ck & \ck \\
                % GraphR                   & 2018 & \ck &     &     &     & \sa &     & \ck &     & \ck & \ck & \ck &     \\
                % SDE                      & 2019 & \ck &     &     &     & \sa & \ck &     & \sc & \ck & \sb & \ck & \sc \\
                ExTensor                 & 2019 & \ck &     &     &     & \sa & \ck & \sb &     & \ck & \sb & \sc & \ck \\
                % PHI                      & 2019 &     & \ck &     & \ck &     &     & \ck & \sc &     & \ck & \ck & \sc \\
                SPiDRE                   & 2019 &     & \ck &     &     & \ck & \ck &     &     & \ck &     & \ck & \sc \\
                ITS                      & 2019 & \ck &     &     &     & \sa &     & \ck &     & \ck & \ck & \ck &     \\
                % SPU                      & 2019 & \ck &     &     &     & \sa & \ck & \ck &     & \ck & \ck & \ck &     \\
                % Tensaurus                & 2020 & \ck &     &     &     & \sa & \sc & \sb &     & \ck & \ck & \ck & \ck \\
                MatRaptor                & 2020 & \ck &     &     &     & \sa &     &     & \ck & \ck & \ck &     & \ck \\
                SpArch                   & 2020 & \ck &     &     &     & \sa &     & \ck &     & \ck & \ck &     & \ck \\
                % Alrescha                 & 2020 & \ck &     &     &     & \sa & \ck &     &     & \ck & \ck & \ck &     \\
                % SIGMA                    & 2020 & \ck &     &     &     & \sa & \ck &     &     & \ck & \ck &     & \ck \\
                % SpaceA                   & 2021 & \ck &     &     &     & \ck & \ck &     &     & \ck & \ck & \ck &     \\
                Gamma                    & 2021 & \ck &     &     &     & \sa &     &     & \ck & \ck & \ck &     & \ck \\
                SPAGHETTI                & 2021 & \ck &     &     &     & \sa &     & \ck &     & \ck & \ck &     & \ck \\
                InnerSP                  & 2021 & \ck &     &     &     & \sa &     &     & \ck & \ck & \ck &     & \ck \\
                % CoSPARSE                 & 2021 & \ck &     &     &     & \sa & \ck & \ck &     & \ck & \ck & \ck &     \\
                % SortCache                & 2021 &     & \ck &     & \ck &     &     & \ck & \sc &     & \ck & \sc & \ck \\
                Sextans                  & 2022 & \ck &     &     &     & \sa &     & \ck &     & \ck & \ck &     & \ck \\
                ASA                      & 2022 &     & \ck & \ck &     &     &     & \sc & \ck &     & \ck & \sc & \ck \\
                % Spada                    & 2023 & \ck &     &     &     & \sa &     &     & \sd & \ck & \ck &     & \ck \\
                Flexagon                 & 2023 & \ck &     &     &     & \sa & \ck & \ck & \ck & \ck & \ck &     & \ck \\
                MOSCON                   & 2023 & \ck &     &     &     & \sa &     & \ck &     & \ck & \ck &     & \ck \\
                GSCSp                    & 2023 & \ck &     &     &     & \sa &     &     & \ck & \ck & \ck &     & \ck \\
                % SDMA                     & 2023 & \ck &     &     &     & \sa &     &     & \sd & \ck & \ck &     & \ck \\
                % SSSR                     & 2023 &     & \ck & \ck &     &     & \ck &     &     & \ck &     & \ck & \sc \\
                % Funnel                   & 2024 & \ck &     &     &     & \sa & \ck &     & \ck & \ck & \ck &     & \ck \\
                GAS                      & 2024 & \ck &     &     &     & \ck &     & \ck &     &     & \ck & \ck & \ck \\
                ACES                     & 2024 & \ck &     &     &     & \sa &     &     & \ck & \ck & \ck &     & \ck \\
                % FullSparse               & 2024 & \ck &     &     &     & \sa &     & \ck &     & \ck & \ck &     & \ck \\
                Sparm                    & 2024 & \ck &     &     &     & \sa & \ck & \ck & \ck & \ck & \ck &     & \ck \\
                % \textbf{SpDCache}~[\ref{publ:SpDC}]           & 2024 &     & \ck &     & \ck &     &     & \ck & \sc &     & \ck & \ck & \sc \\
                SPADES                   & 2024 & \ck &     &     &     & \sa & \sc & \sc & \sc & \ck &     & \ck & \sc \\
                % DySpMM                   & 2024 & \ck &     &     &     & \sa & \ck &     &     & \ck & \ck &     & \ck \\
                Vesper                   & 2024 & \ck &     &     &     & \sa & \ck & \ck & \ck & \ck & \ck & \ck & \ck \\
                % Mentor                   & 2024 & \ck &     &     &     & \sa &     &     & \ck & \ck & \ck &     & \ck \\
                IOPS                     & 2025 & \ck &     &     &     & \sa & \ck & \ck &     & \ck & \ck &     & \ck \\
                % Leda                     & 2025 & \ck &     &     &     & \sa &     & \ck &     & \ck & \ck &     & \ck \\
                % DMSA                     & 2025 & \ck &     &     &     & \sa &     &     & \sd & \ck & \ck &     & \ck \\
                C\textsuperscript{3}ache & 2025 &     & \ck &     & \ck &     &     & \sd & \sd & \ck & \ck &     & \ck \\
                ScanNow                  & 2025 & \ck &     &     &     & \sa &     &     & \sd & \ck & \ck &     & \ck \\
                %
                % \sa \; Most standalone accelerators connect directly to main memory. \\
                % \sb \; Not main contribution but addressed in the article. \\
                % \sc \; Could be adapted but not discussed in the article. \\
                % \sd \; New dataflow based on one or several algorithms.
            \end{tblr}
        }

    \end{minipage}
    \hfill
    \begin{minipage}[t]{0.30\linewidth}
        \raggedright

        \smallskip
        
        \pause[4]%
        \arrow[\bfseries\color{Red}]{Diversité} des choix de conception, toutes les caractéristiques sont couvertes
        %
        \pause[2]%
        \begin{tikzpicture}[remember picture, overlay]
            \node [shift={(2.4cm, -1.65cm)}, single arrow, very thick, fill=Red, opacity=0.3, minimum width = 1cm, single arrow head extend=3pt, minimum height=4.85cm, rotate=-90] at (current page.west)
                {}; % length of arrow
        \end{tikzpicture}%
        %
        \pause[3]%
        {%
            \def\X{3.39cm}%
            \def\Y{1.60cm}%
            \def\H{38pt}%
            \def\Wa{0.88cm}%
            \def\Wb{1.52cm}%
            \begin{tikzpicture}[remember picture, overlay]
                \node [shift={(\X, \Y)}, rotate=0] at (current page.west)
                    {%
                        \lbox[standard jigsaw,opacityframe=1.,opacityback=0.]{Cyan}{%
                            \rule{0pt}{\H}%
                            \hspace*{\Wa}%
                        }%
                    };
            \end{tikzpicture}%
            %
            \begin{tikzpicture}[remember picture, overlay]
                \node [shift={(\X + 1.67cm, \Y)}, rotate=0] at (current page.west)
                    {%
                        \lbox[standard jigsaw,opacityframe=1.,opacityback=0.]{Purple}{%
                            \rule{0pt}{\H}%
                            \hspace*{\Wb}%
                        }%
                    };
            \end{tikzpicture}%
            %
            \begin{tikzpicture}[remember picture, overlay]
                \node [shift={(\X + 3.66cm, \Y)}, rotate=0] at (current page.west)
                    {%
                        \lbox[standard jigsaw,opacityframe=1.,opacityback=0.]{Green}{%
                            \rule{0pt}{\H}%
                            \hspace*{\Wb}%
                        }%
                    };
            \end{tikzpicture}%
            %
            \begin{tikzpicture}[remember picture, overlay]
                \node [shift={(\X + 5.32cm, \Y)}, rotate=0] at (current page.west)
                    {%
                        \lbox[standard jigsaw,opacityframe=1.,opacityback=0.]{Blue}{%
                            \rule{0pt}{\H}%
                            \hspace*{\Wa}%
                        }%
                    };
            \end{tikzpicture}%
            %
            \begin{tikzpicture}[remember picture, overlay]
                \node [shift={(\X + 6.65cm, \Y)}, rotate=0] at (current page.west)
                    {%
                        \lbox[standard jigsaw,opacityframe=1.,opacityback=0.]{Orange}{%
                            \rule{0pt}{\H}%
                            \hspace*{\Wa}%
                        }%
                    };
            \end{tikzpicture}%
            %
            \pause[1]%
            \begin{tikzpicture}[remember picture, overlay]
                \node [shift={(\X - 1.8cm, \Y + 0.5cm)}, rotate=0] at (current page.west)
                    {%
                        \begin{lightbox}[left=0pt,right=0pt,top=0pt,bottom=-1pt,colback=black!5]{black!10}%
                            \itshape\scriptsize Contribution \cref{ref:HiPEAC}%
                        \end{lightbox}%
                    };
            \end{tikzpicture}%
        }

        \pause[4]%
        \vspace*{0.75cm}

        \qquad\makebox(0,0){\put(0,0.0\normalbaselineskip){%
            $\left.\rule{0pt}{2.7\normalbaselineskip}\right\}$
            \parbox[t]{0.9\linewidth}{
                \raggedright
                Exploration des algorithmes \textbf{\textcolor{Red}{IP et OP}}
            }
        }}

        \vspace*{1.1cm}

        \quad\makebox(0,0){\put(0,0.0\normalbaselineskip){%
            $\left.\rule{0pt}{3.3\normalbaselineskip}\right\}$
            \parbox[t]{0.9\linewidth}{
                \raggedright
                Exploration de l'algorithme de \textbf{\textcolor{Red}{Gustavson}}
            }
        }}

        \vspace*{1.2cm}

        \makebox(0,0){\put(0,0.0\normalbaselineskip){%
            $\left.\rule{0pt}{3.0\normalbaselineskip}\right\}$
            \parbox[t]{0.9\linewidth}{
                \raggedright
                \textbf{\textcolor{Red}{Combinaisons d'algorithmes}}
            }
        }}

    \end{minipage}

\end{frame}

\refStep{ref:HiPEAC}

\begin{frame}
    \frametitle{Choix de conception}

    \vspace*{-7pt}
    \def\H{6.5cm}

    \begin{minipage}[t][\H]{0.44\linewidth}
        \raggedright

        \pause[1]%
        \textbf{\textcolor{Mulberry}{Comment choisir une architecture ?}}

        \smallskip

        \pause[2]%
        \point[\bfseries\color{Red}][0]{Accélérateur autonome}
        
        \point[\bfseries\color{Red}][0]{spécifique, performance}[=]

        \smallskip
        
        \arrow[][2]{\textit{Kernel} spécifique}[\faCheckCircleO]
        
        \arrow[][2]{Application spécifique (parfois)}[\faCheckCircleO]
        
        \arrow[][2]{Mémoires customisées}[\faCheckCircleO]
        
        \arrow[][2][Green]{Hautes performances}[\faPlus]
        
        \arrow[][2][Red]{Surface conséquente}[\faMinus]

        \smallskip

        \pause[3]%
        \point[\bfseries\color{Red}][0]{Assisté par processeur}
        
        \point[\bfseries\color{Red}][0]{flexible, faible coût}[=]

        \smallskip
        
        \arrow[][2][Green]{Pour tous \textit{kernels}}[\faPlus]
        
        \arrow[][2]{Mémoires existantes}[\faCheckCircleO]
        
        \arrow[][2][Green]{Intégrable dans divers systèmes}[\faPlus]
        
        \arrow[][2][Red]{Intégration peu évidente}[\faMinus]
        
        \arrow[][2][Red]{Performances moins élevées}[\faMinus]
        
        \arrow[][2][Green]{Faible surface ajoutée}[\faPlus]

    \end{minipage}
    \hfill
    \begin{minipage}[t][\H]{0.54\linewidth}
        \centering

        \pause[1]%
        \includegraphics<1>[width=\linewidth]{Figures/PDF/SoA_AnalysisForDesigners_0.pdf}
        \includegraphics<2>[width=\linewidth]{Figures/PDF/SoA_AnalysisForDesigners_1.pdf}
        \includegraphics<3>[width=\linewidth]{Figures/PDF/SoA_AnalysisForDesigners_2.pdf}

    \end{minipage}

    \pause[1]%
    \vspace{0.2cm}
    \drefX{ref:HiPEAC}{Valentin Isaac-{-}Chassande, Adrian Evans, Yves Durand, and Frédéric Rousseau. 2024. Dedicated Hardware Accelerators for Processing of Sparse Matrices and Vectors: A Survey. ACM Trans. Archit. Code Optim. 21, 2, Article 27 (June 2024), 26 pages.
    %
    \begin{tikzpicture}[remember picture, overlay]
        \node [shift={(-3.8cm, 2.80cm)}, rotate=0] at (current page.east)
            {%
                {\itshape\scriptsize Contribution \cref{ref:HiPEAC}}%
            };
    \end{tikzpicture}%
    }

\end{frame}

\subsection{Direction de la thèse}

\begin{frame}
    \frametitle{Quelle est la meilleure solution ?}
    
    Les accélérateurs matériels dédiés (autonomes et assistés)

    \point[\bfseries][2][Green]{Répondent efficacement à la problématique du calcul creux}[\ding{51}]

    \bigskip

    \point[][2][Red]{Méthodes et choix de conception \textbf{divers}}
    \arrow[]{\underbar{Comparaison de performances exhaustive difficile}}

    \smallskip

    \point[][2][white]{Ce qui resort :}
    \arrow[\bfseries][0][Red]{Pour les \textit{kernels} Matrice-Matrice : \textcolor{Red}{Gustavson}}
    
    \subpt[][35][Red]{Réduction théorique du trafic par rapport à IP et OP} \cref{ref:HiPEAC}

    \smallskip

    \point[\color{white}][2]{Ce qui ressort :}
    \arrow[\bfseries][0][Red]{Pour le SpMV : \textcolor{Red}{Outer-Product}}

    \subpt[][35][Red]{Avantages théoriques d'optimisation par rapport à IP} \cref{ref:HiPEAC}

    \smallskip

    \arrow[\bfseries\color{Red}][4]{Gust et OP :} Algorithmes \underbar{les plus utilisés} dans les accélérateurs \underbar{les plus performants}

    \medskip

    \arrow[\bfseries\color{Mulberry}]{Besoin de support pour les réductions \og aléatoires \fg{}}

    \bigskip

    En général, pour tous les accélérateurs :

    \arrow[][2][Red]{Grandes \textbf{variations de performances} selon les matrices et leurs \textbf{structures}}

    \medskip

    \arrow[\bfseries\color{Mulberry}]{Besoin de considérer les structures des matrices}

\end{frame}

\begin{frame}
    \frametitle{Direction de cette thèse}

    \textbf{\textcolor{Red}{Hypothèses :}}

    \smallskip

    \point[][0][Red]{Noyau de calcul :} \textbf{SpMV} (multiplication matrice creuse-vecteur)

    \smallskip

    \point[][0][Red]{Algorithme :} \textbf{OP} (Outer-Product)

    \smallskip
    
    \point[][0][Red]{Problème :} Scatter \arrow[\bfseries]{Réductions \og aléatoires \fg{}}

    \bigskip

    \point[][0][Red]{Architecture de \textbf{type générique}, \textbf{assisté par processeur}, faible coût en surface}

    \arrow[][2][Red]{Au niveau des \textbf{caches de données}}

    \bigskip

    \textbf{\textcolor{Mulberry}{Objectif 1 :} Réduire l'impact de l'irrégularité des réductions en analysant le comportement des caches de données}

    \medskip

    \textbf{\textcolor{Mulberry}{Objectif 2 :} Comprendre le rôle des structures des matrices dans les performances}

    \bigskip

    \textbf{\textcolor{Red}{Approche :}}

    \smallskip
    
    \arrow[\bfseries][0][Red]{SpDCache :} Cache de données spécialisé qui adresse les réductions \og aléatoires \fg{}

\end{frame}